\section{Lifted finite cocycles at a split column}

% BEGIN F82L1
\paragraph{F82L1. The two evaluation maps over the ordered roots.}
Use the twisted split middle-range index and good point $\pi$ of
G81S1. Use the integral normalization, good coefficient enlargement
$L/\mathbb Q_p$ with integers $O$, and ordered-root cover $Y\to X$
of G81S3. The enlargement can include the fields needed for all
labelled local comparisons, since there are finitely many places.
Write $y_0$ for the distinguished ordering above $\pi$, $A$ for
the algebra spanned by the family matrices over the ordered-root
Hecke algebra, and $U(g)=\chi_1(g)^{-1}R(g)$.
The evaluation on the rank summand at $y_0$ is
$\psi_\rho:A\to M_r(O)$, and the evaluation on the first
character is $\psi_1:A\to O$. D79G1--D79G3 give an integral
idempotent $e$ with $\psi_\rho(e)=E_{11}$ and a surjection onto
the full rank matrix algebra. Denote by $E$ the integral polynomial
in the character corner of G81S3. Its sole nonzero block at $y_0$
is $\psi_1(E)=\delta$, where $v_p(\delta)=v_\delta(i)$.

Give the left ideals $Ae$ and $AE$ the local action by left
multiplication by $U(g)$. Put $T=O^r$ with action
$W=\rho_i(\pi)/\chi_1$, and write $W$ also for $T[1/p]$.
There are equivariant integral surjections
\[
 Ae\longrightarrow T,\quad m\longmapsto\psi_\rho(m){\bf e}_1,
 \qquad
 AE\longrightarrow\delta O,\quad m\longmapsto\psi_1(m),
\]
where the action on $\delta O$ is trivial. For the first, prescribe
the first column of $\psi_\rho(a)$ and take $m=ae$.
For the second, $\psi_1(aE)=\delta\psi_1(a)$, and scalars
realize every multiple of $\delta$. These maps are $A$-linear for
the indicated block actions. Each left ideal is a submodule of
the finite product of integral matrices and hence finite free
over $O$. After rationalization, mapping each matrix at each
$y$ to its list of columns embeds either ideal equivariantly into
the direct sum of $n_V$ copies of $R_{p(y)}/\chi_1$ for each $y$,
where $p(y)$ denotes the point below $y$. Both rational targets
are quotients of subobjects of this sum. Consequently these exact
surjections are two tests in L81S4 at every column $u$, with
the same $a=X.e_S(i)$. In particular the second target uses its
actual image lattice $\delta O$, even for nonunit $\delta$.
% END F82L1

% BEGIN F82L2
\paragraph{F82L2. An integral finite basis and the degree character.}
Fix $u$, put $K=E_u$, $G=G_K$, and identify $G$ with the
specified decomposition subgroup. For a stable lattice $P$, write
$H^1_f(K,P)$ for the inverse image of $H^1_f(K,P[1/p])$ in
continuous integral cohomology. The subgroup is saturated in
$H^1(K,P)$, since its defining rational quotient is a vector
space. A81R1 gives finite generation of integral cohomology.
Every continuous rational cocycle on the compact group has bounded
image and hence admits a common nonzero integral multiple with
values in $P$. A rational coboundary between integral cocycles
becomes an integral coboundary after multiplying its vector by
a denominator. These two observations give
$H^1_f(K,P)[1/p]=H^1_f(K,P[1/p])$.

For $P=T$ let $j=\dim_L H^1_f(K,W)\le [K:\mathbb Q_p]m$
as in G81S1. Over the DVR the quotient of $H^1_f(K,T)$ by its
torsion has an $O$-basis of $j$ classes. Lift its basis elements
to integral finite cocycles $\kappa_{u,1},\ldots,\kappa_{u,j}$.
The difference between any integral finite class and a suitable
integral combination is torsion. An integral combination that
is torsion has all coefficients zero. The kernel of integral to
rational cohomology is the image of $(W/T)^G$, by the invariants
sequence. Since $W^G=0$ by G81S1, that image is isomorphic to
$(W/T)^G$ and is precisely the torsion of the finite subgroup.
Its length divided by $\operatorname{length}_O(O/pO)$ is the
column term $h_W(i,u,\pi)$ on these coefficients.
For $b_s\in L/O$, the cocycle $\sum_s b_s\kappa_{u,s}$
in $W/T$ is finite in the propagated condition: lifting the $b_s$
to $L$ gives a rational finite linear combination. A different
choice of lifts gives the same function modulo $T$. The argument
also covers $j=0$.

Let $d_u:G\to\mathbb Z_p\subset O$ be reduction to the
residue-field Galois group followed by its $p$-component,
using arithmetic Frobenius as the generator. Then $d_u$ vanishes
on inertia and equals $1$ on a Frobenius lift. The ideal-valued
cocycle $\delta d_u$ is finite, since its rational multiple of
the unramified degree line is crystalline by J80H1. If global
functions are required as in the column condition, extend each
$\kappa_{u,s}$ and $d_u$ by zero outside the decomposition
subgroup. Only their restrictions to that subgroup enter the
cocycle and finiteness tests.
% END F82L2

% BEGIN F82L3
\paragraph{F82L3. Lifting the augmented columns with the stored exponent.}
For either map $f:P\to Q$ in F82L1, the split lifting conclusion
L81S4 gives $p^aH^1_f(K,Q)\subseteq f_*H^1_f(K,P)$.
This holds over the particular good coefficient ring of G81S3
and for every integral finite target class, including torsion.
It gives equality of functions after choosing representatives.
Indeed, if $z$ is the target cocycle and
$f(\widehat z)-p^az$ is the coboundary of $q\in Q$, subtract
the coboundary of any lift of $q$ in $P$ from $\widehat z$.
The resulting source cocycle is still continuous and finite.
Thus there are integral finite functions
\[
 z_s:G\to Ae,\qquad z_0:G\to AE
\]
whose evaluations are exactly $p^a\kappa_{u,s}$ and
$p^a\delta d_u$, respectively. Each source is itself a lattice,
so these functions give the finite integral cohomology classes
defined through rationalization of the left ideal. Conversely any
abstract cocycle representative of such a class differs from a
finite continuous representative by an integral coboundary.
The latter is continuous because the action is continuous.
This identifies the finite classes used by the two lift tests
with the integral finite classes just constructed.
The factor $p^a$ works simultaneously on all linear combinations.
In the scalar-valued matrix evaluation of the last column the
factor is $p^a\delta$, of normalized valuation $a+v_\delta(i)$.
% END F82L3

% BEGIN F82L4
\paragraph{F82L4. The finite cocycle space in every evaluated fibre.}
At $y\in Y$ write $V_y=L^{n_V}$ with local action given by
evaluating $U(g)$ at $y$. Put
\[
  \mathcal L_{y,u}=Z_f^1(K,V_y)
     =\{z\in Z^1_{\rm cont}(K,V_y):
                 [z]\mapsto0\text{ in }H^1(K,B_{\rm cris}\otimes_{\mathbb Q_p}V_y)\}.
\]
This is an $L$-linear subspace of all functions $G\to V_y$.
It contains the image of $c_y:v\mapsto(g\mapsto(U_y(g)-1)v)$.
For $P=Ae$ or $AE$ and any $A$-linear $F:P\to V_y$, scalars
make $F$ $O$-linear. It therefore extends uniquely to a continuous
$L$-linear map $P[1/p]\to V_y$ by division by scalars. It
intertwines $G$ because $U(g)\in A$. Extending it to
$B_{\rm cris}$ sends rational finite cocycles to finite cocycles.
F82L3's representative observation now shows
$g\mapsto F(z(g))\in\mathcal L_{y,u}$ for every integral
representative whose class in either left ideal is finite.
This proves both universal image conditions of the original
column space, including maps other than the coordinate evaluations.

The kernel of $c_y$ is $V_y^G$, and the quotient of
$\mathcal L_{y,u}$ by its coboundary subspace is precisely
$H^1_f(K,V_y)$. Each $V_y$ is crystalline by L81S1. By S80J1
its negative jump multiplicities at each embedding are those of
$W$: the difference of the full and rank lists is the character
string $0,\ldots,t-1$. The crystalline dimension formula J80H1
and $W^G=0$ give
\[
 \dim_L H^1_f(K,V_y)
 =\dim_L V_y^G+\dim_L\bigl(D_{\rm dR}(W)/\mathrm{Fil}^0\bigr)
 =\dim_L\ker c_y+j .
\]
Here the de Rham dimension counts negative jumps at all embeddings
of $K$ in $L$. Consequently $\dim_L\mathcal L_{y,u}=n_V+j$,
even when the invariant dimension of the fibre changes.
In particular the required quotient rank bound holds for every
$y$ above the complement of the specified Hecke support. These
spaces and F82L1--F82L3 provide the local augmented-column data
on one ordered-root cover with the prescribed lifting exponent.
% END F82L4
